\section{Введение}
\label{sec:Chapter0} \index{Chapter0}

\subsection{Актуальность темы}

Речевой сигнал является одним из основных каналов передачи информации между
человеком и вычислительными системами. Современные приложения --- голосовые
помощники, системы автоматического распознавания речи (ASR), синтеза речи по
тексту (TTS), телеконференцсвязь, слуховые аппараты --- предъявляют высокие
требования к чистоте и разборчивости речи. Однако в реальных условиях речевой
сигнал всегда искажён: к нему примешиваются аддитивный шум окружения,
реверберация помещения, искажения микрофонного тракта и артефакты сжатия
кодеками. Эти искажения снижают как субъективное качество восприятия, так и
точность работы последующих алгоритмов обработки речи.

Задача восстановления речи (speech restoration, SR) состоит в преобразовании
искажённого сигнала в сигнал студийного качества. В отличие от классического
улучшения речи (speech enhancement, SE), которое преимущественно нацелено на
подавление аддитивного шума, восстановление речи рассматривает широкий спектр
деградаций одновременно и стремится не просто <<вычесть>> помеху, а
синтезировать чистый сигнал заново. Такой подход особенно востребован при
подготовке обучающих корпусов для генеративных речевых моделей: качество
синтезируемой речи напрямую зависит от качества данных, а ручной сбор
студийных записей в больших объёмах экономически нецелесообразен. Так, при
помощи модели Miipher~\cite{miipher1} многоговорящий корпус LibriTTS был
очищен и преобразован в корпус студийного качества LibriTTS-R~\cite{librittsr},
что заметно повысило качество обученных на нём систем синтеза.

Прорыв в области обеспечило применение представлений, полученных методами
самоконтролируемого обучения (self-supervised learning, SSL). Модели wav2vec
2.0~\cite{wav2vec2}, HuBERT~\cite{hubert} и WavLM~\cite{wavlm}, предобученные
на тысячах часов немаркированной речи, формируют устойчивые к шуму
лингвистические и акустические представления. Архитектура
Miipher-2~\cite{miipher2} развивает эту идею: вместо обработки сигнала во
временной или частотной области очистка выполняется в пространстве признаков
SSL-энкодера, после чего нейросетевой вокодер синтезирует чистую волновую
форму. Подобная схема демонстрирует устойчивость к разнообразным деградациям
и масштабируется на обработку данных объёмом в миллионы часов.

\subsection{Проблема и мотивация работы}

Несмотря на высокое качество, оригинальная модель Miipher-2 опирается на
закрытый энкодер USM~\cite{usm} с двумя миллиардами параметров, предобученный
на 12~миллионах часов аудио. Ни веса USM, ни обученный на их основе очиститель
признаков не опубликованы и недоступны для воспроизведения. Это делает
невозможным прямое использование Miipher-2 в открытых проектах и научных
исследованиях и ставит вопрос: можно ли заменить закрытый энкодер открытыми
SSL-моделями, сохранив приемлемое качество восстановления речи?

Открытые энкодеры WavLM и HuBERT обучены на существенно меньших объёмах данных
и имеют другую структуру представлений. При этом известно, что разные слои
SSL-моделей кодируют разную информацию: нижние слои ближе к акустике, верхние
--- к фонетике и семантике~\cite{superb}. Следовательно, выбор слоя извлечения
признаков и типа вокодера становится самостоятельной исследовательской задачей,
от решения которой напрямую зависит итоговое качество системы.

\subsection{Объект и предмет исследования}

\textbf{Объект исследования} --- нейросетевые системы восстановления и очистки
речевых сигналов на основе самоконтролируемых акустических представлений.

\textbf{Предмет исследования} --- влияние выбора открытого SSL-энкодера, номера
слоя извлечения признаков и типа нейросетевого вокодера на качество и
вычислительную эффективность системы восстановления речи, построенной по схеме
Miipher-2.

\subsection{Цель и задачи}

\textbf{Цель работы} --- разработать и исследовать систему восстановления речи
на базе архитектуры Miipher-2 с заменой закрытого энкодера USM на открытые
SSL-модели и определить конфигурацию, обеспечивающую наилучший баланс качества
и быстродействия.

Для достижения цели поставлены следующие \textbf{задачи}:
\begin{enumerate}
    \item провести аналитический обзор классических и нейросетевых методов
        улучшения и восстановления речи, а также SSL-энкодеров и нейросетевых
        вокодеров;
    \item формализовать задачу восстановления речи и сформулировать требования
        к разрабатываемой системе;
    \item воспроизвести архитектуру Miipher-2 с открытыми компонентами:
        энкодерами WavLM и HuBERT, вокодерами HiFi-GAN и WaveFit;
    \item исследовать зависимость качества восстановления от номера слоя
        извлечения признаков SSL-энкодера;
    \item экспериментально сравнить четыре комбинации <<энкодер — вокодер>> по
        объективным метрикам качества и вычислительной эффективности;
    \item выработать рекомендации по выбору конфигурации системы.
\end{enumerate}

\subsection{Методы исследования}

В работе используются методы цифровой обработки сигналов (кратковременное
преобразование Фурье, мел-спектральный анализ), аппарат глубокого обучения
(свёрточные и трансформерные нейронные сети, генеративно-состязательное
обучение), методы трансферного обучения с замороженными предобученными
энкодерами, а также методики объективной оценки качества речи (PESQ, STOI,
SI-SDR, DNSMOS) и вычислительный эксперимент.

\subsection{Научно-теоретическая и практическая значимость}

Полученная система позволяет очищать зашумлённые речевые записи до качества,
близкого к студийному, опираясь исключительно на открытые компоненты. Это даёт
возможность автоматически формировать высококачественные обучающие корпуса для
систем синтеза и распознавания речи, восстанавливать архивные и пользовательские
аудиозаписи, а также повышать разборчивость речи в телекоммуникационных
приложениях. Выработанные рекомендации по выбору слоя извлечения признаков и
типа вокодера применимы к широкому классу систем, построенных на
SSL-представлениях, и обладают самостоятельной методической ценностью.

\subsection{Структура работы}

Работа состоит из введения, четырёх основных разделов, заключения, списка
использованных источников и приложения. В разделе~\ref{sec:Chapter1}
формализуется задача и формулируются требования к решению.
Раздел~\ref{sec:Chapter2} содержит обзор существующих подходов и оценку их
соответствия требованиям. В разделе~\ref{sec:Chapter3} описывается декомпозиция
задачи и предлагаемая архитектура. Раздел~\ref{sec:Chapter4} посвящён
программной реализации и результатам экспериментов. В заключении подводятся
итоги исследования и формулируются рекомендации.
